Un satellite d'observation en orbite à \qty{\v{H}}{km} d'altitude doit pouvoir distinguer au sol deux objets séparés de \qty{\v{s}}{cm}. On considère une lumière de longueur d'onde \qty{\v{lam}}{nm} et on néglige les effets de l'atmosphère.

\begin{enumerate}[label=\alph*)]
    \item Quel angle le télescope du satellite doit-il pouvoir résoudre ?
    \item Quel diamètre minimal doit avoir son miroir ?
    \item Le miroir du télescope spatial Hubble mesure \qty{2.4}{m} de diamètre. À la même altitude, quel est le plus petit détail qu'un tel miroir pourrait distinguer au sol ?
\end{enumerate}
